\SourceLeaf{022}Ut hic error medius fiat quam minimus, statuere oportebit
\[
\begin{aligned}
\alpha'&=\frac{n^4-m^4}{n'{}^4-m'{}^4}.\mu'\mu'\\[4pt]
\alpha''&=\frac{n^4-m^4}{n''{}^4-m''{}^4}.\mu''\mu''\text{ etc.}
\end{aligned}
\]
Manifesto hi valores evolvi nequeunt, nisi insuper ratio quantitatum \(n,n',n''\) etc. ad \(m,m',m''\) etc. aliunde nota fuerit: qua cognitione exacta deficiente, saltem tutissimum videtur\footnote{Scilicet cognitionem quantitatum \(\mu',\mu''\) etc. in eo solo casu in potestate esse concipimus, ubi per rei naturam errores \(x,x',x''\) etc. ipsis \(1,\mu',\mu''\) etc. proportionales, aeque probabiles censendi sunt, aut potius ubi
\[
\varphi x=\mu'\varphi'(\mu'x)=\mu''\varphi''(\mu''x)\text{ etc.}
\]
}, illas his proportionales supponere (v. art. 11), unde prodeunt valores
\[
\alpha'=\frac1{\mu'\mu'},\quad\alpha''=\frac1{\mu''\mu''}\text{ etc.}
\]
i. e. coëfficientes \(\alpha',\alpha''\) etc. aequales statui debent ponderibus relativis observationum, ad quas pertinent errores \(x',x''\) etc., assumto pondere observationis, ad quam pertinet error \(x\), pro unitate. Hoc pacto, designante ut supra \(\sigma\) multitudinem errorum propositorum, habebimus valorem medium expressionis
\[
\frac{xx+\alpha'x'x'+\alpha''x''x''+\text{etc.}}{\sigma}
\]
\(=mm\), atque errorem medium metuendum, dum valorem fortuito determinatum huius expressionis pro valore vero ipsius \(mm\) adoptamus
\[
\frac{\BareRoot{(n^4+\alpha'\alpha'n'{}^4+\alpha''\alpha''n''{}^4+\text{etc.}-\sigma m^4)}}{\sigma}
\]
et proin, siquidem licet, ipsas \(n,n',n''\) etc. ipsis \(m,m',m''\) proportionales supponere,
\[
\mathord{\cdot}=\BareRoot{\frac{n^4-m^4}{\sigma}}
\]
quae formula identica est cum ea, quam supra pro casu observationum eiusdem generis inveneramus.

\Article{17}
Si valor quantitatis, quae ab alia quantitate incognita pendet, per observa\SourceLeaf{023}tionem praecisione absoluta non gaudentem determinata est, valor incognitae hinc calculatus etiam errori obnoxius erit, sed nihil in hac determinatione arbitrio relinquitur. At si \textit{plures} quantitates ab eadem incognita pendentes per observationes haud absolute exactas innotuerunt, valorem incognitae vel per quamlibet harum observationum eruere possumus, vel per aliquam plurium observationum combinationem, quod infinitis modis diversis fieri potest. Quamquam vero valor incognitae tali modo prodiens errori semper obnoxius manet, tamen in alia combinatione maior, in alia minor error metuendus erit. Similiter res se habebit, si plures quantitates a pluribus incognitis simul pendentes sunt observatae: prout observationum multitudo multitudini incognitarum vel aequalis, vel hac minor vel maior fuerit, problema vel determinatum, vel indeterminatum, vel plus quam determinatum erit (generaliter saltem loquendo), et in casu tertio ad incognitarum determinationem observationes infinitis modis diversis combinari poterunt. E tali combinationum varietate eas eligere, quae maxime ad rem faciant, i. e. quae incognitarum valores erroribus minimis obnoxios suppeditent, problema sane est in applicatione matheseos ad philosophiam naturalem longe gravissimum.

In Theoria motus corporum coelestium ostendimus, quomodo valores incognitarum \textit{maxime probabiles} eruendi sint, si lex probabilitatis errorum observationum cognita sit; et quum haec lex natura sua in omnibus fere casibus hypothetica maneat, theoriam illam ad legem maxime plausibilem applicavimus, ubi probabilitas erroris \(x\) quantitati exponentiali \(e^{-hhxx}\) proportionalis supponitur, unde methodus a nobis dudum in calculis praesertim astronomicis, et nunc quidem a plerisque calculatoribus sub nomine methodi quadratorum minimorum usitata demanavit.

Postea ill. \textsc{Laplace}, rem alio modo aggressus, idem principium omnibus aliis etiamnum praeferendum esse docuit, quaecunque fuerit lex probabilitatis errorum, si modo observationum multitudo sit permagna. At pro multitudine observationum modica, res intacta mansit, ita ut si lex nostra hypothetica respuatur, methodus quadratorum minimorum eo tantum nomine prae aliis commendabilis habenda sit, quod calculorum concinnitati maxime est adaptata.

Geometris itaque gratum fore speramus, si in hac nova argumenti tractatione docuerimus, methodum quadratorum minimorum exhibere combinationem ex omnibus optimam. non quidem proxime, sed absolute, quaecunque fuerit lex probabilitatis errorum, quaecunque observationum multitudo, si modo notionem \SourceLeaf{024}erroris medii non ad mentem ill. \textsc{Laplace}, sed ita, ut in artt. 5 et 6 a nobis factum est, stabiliamus.

Ceterum expressis verbis hic praemonere convenit, in omnibus disquisitionibus sequentibus tantummodo de erroribus irregularibus atque a parte constante liberis sermonem esse, quum proprie ad perfectam artem observandi pertineat, omnes errorum constantium caussas summo studio amovere. Quaenam vero subsidia calculator tales observationes tractare suscipiens, quas ab erroribus constantibus non liberas esse iusta suspicio adest, ex ipso calculo probabilium petere possit, disquisitioni peculiari alia occasione promulgandae reservamus.

\Article{18}
\textsc{Problema.} \textit{Designante \(U\) functionem datam quantitatum incognitarum \(V,V',V''\) etc., quaeritur error medius \(M\) in determinatione valoris ipsius \(U\) metuendus, si pro \(V,V',V''\) etc. adoptentur non valores veri, sed ii, qui ex observationibus ab invicem independentibus, erroribus mediis \(m,m',m''\) etc. resp. obnoxiis prodeunt.}

\textit{Sol.} Denotatis erroribus in valoribus observatis ipsarum \(V,V',V''\) etc. per \(e,e',e''\) etc., error inde redundans in valorem ipsius \(U\) exprimi poterit per functionem linearem
\[
\lambda e+\lambda'e'+\lambda''e''+\text{etc.}=E
\]
ubi \(\lambda,\lambda',\lambda''\) etc. sunt valores quotientium differentialium \(\frac{\mathrm{d}U}{\mathrm{d}V},\frac{\mathrm{d}U}{\mathrm{d}V'},\frac{\mathrm{d}U}{\mathrm{d}V''}\) etc. pro valoribus veris ipsarum \(V,V',V''\) etc. siquidem observationes satis exactae sunt, ut errorum quadrata productaque negligere liceat. Hinc primo sequitur, quoniam observationum errores a partibus constantibus liberi supponuntur, valorem medium ipsius \(E\) esse \(=0\). Porro error medius in valore ipsius \(U\) metuendus erit radix quadrata e valore medio ipsius \(EE\), sive \(MM\) erit valor medius aggregati
\[
\begin{aligned}
\lambda\lambda ee+\lambda'\lambda'e'e'+\lambda''\lambda''e''e''+\text{etc.}{}\\
{}+2\lambda\lambda'ee'+2\lambda\lambda''ee''+2\lambda'\lambda''e'e''+\text{etc.}
\end{aligned}
\]
At valor medius ipsius \(\lambda\lambda ee\) fit \(\lambda\lambda mm\), valor medius ipsius \(\lambda'\lambda'e'e'\) fit \(=\lambda'\lambda'm'm'\) etc.; denique valores medii productorum \(2\lambda\lambda'ee'\) etc. omnes fiunt \(=0\). Hinc itaque colligimus
\[
M=\BareRoot{(\lambda\lambda mm+\lambda'\lambda'm'm'+\lambda''\lambda''m''m''+\text{etc.})}
\]
